\section{Introduccion}

La discusion sobre salario minimo suele enfocarse en empleo y salarios, pero una pregunta central para macro y politica social es el posible traspaso hacia precios (``espiral salario--precios''). En economias con alta heterogeneidad laboral e informalidad, la magnitud de ese traspaso puede variar entre periodos y canastas de consumo.

La literatura internacional documenta que el salario minimo puede generar ajustes de precios en rubros intensivos en trabajo de bajos salarios \citep{aaronson2008, cengiz2019}. Tambien muestra que los efectos distributivos y de formalizacion dependen del contexto institucional y de enforcement \citep{card1995, dube2019}. En este marco, Paraguay ofrece un caso relevante por la combinacion de reajustes discretos del SML, dinamica inflacionaria y cambios recientes en indicadores de hogares.

Este trabajo contribuye con una plataforma reproducible y un articulo tecnico centrado en cuatro objetivos:
\begin{enumerate}
  \item cuantificar diferencias pre/post de inflacion en ventanas alrededor de ajustes del SML;
  \item comparar esas diferencias con distribuciones placebo (Monte Carlo);
  \item integrar contexto laboral (R02) y de hogares (R01) para lectura economica de los hallazgos;
  \item documentar limites de interpretacion para uso responsable de evidencia en decision publica.
\end{enumerate}

El enfoque es deliberadamente transparente: todos los cuadros y figuras del articulo se generan desde scripts reproducibles incluidos en esta carpeta.

Como aporte aplicado para divulgacion y uso por tomadores de decision, el estudio se complementa con una app web interactiva del proyecto, disponible en \url{https://diegomezapy.github.io/estudioSMLparaguay/}.
